\chapter{Instruction Formats}
\label{ch:instruction-formats}

\section{Overview}

All SIRCIS instructions are exactly 32 bits (4 bytes, or 2 words) in length. This fixed instruction size simplifies
instruction fetch and decode logic, though it may result in larger code size compared to variable-length instruction
sets.

The SIRC-1 supports three distinct instruction formats:

\begin{enumerate}
	\item Immediate
	\item Short Immediate (with Shift)
	\item Register
\end{enumerate}

\textbf{Note:} Operandless assembler meta-instructions such as \mnemonic{NOOP} still assemble to one of these three CPU
instruction formats.

Each format is optimized for different types of operations and addressing modes.

\section{Format Overview}

\begin{table}[H]
	\centering
	\small
	\begin{tabular}{lp{9cm}}
		\toprule
		\textbf{Format} & \textbf{Typical Use}                                  \\
		\midrule
		Immediate       & Operations with a 16-bit constant value               \\
		Short Immediate & Operations with an 8-bit constant with optional shift \\
		Register        & Operations on multiple registers with optional shift  \\
		\bottomrule
	\end{tabular}
	\caption{Instruction Format Summary}
	\label{tab:format-summary}
\end{table}

\section{Common Fields}

All instruction formats share some common fields:

\begin{description}
	\item[Opcode (6 bits)] Identifies the operation to perform. With 6 bits, up to 64 distinct operations can be encoded.

	\item[Condition Flags (4 bits)] Specifies the conditions under which the instruction executes. See
	      Chapter~\ref{ch:condition-codes} for details.

	\item[Additional Flags (2 bits)] Used in memory operations to specify address register pairs (\reg{l}, \reg{a}, \reg{s},
	      \reg{p}) and in ALU operations to specify how the status register is updated.

\end{description}

\section{Immediate Format}

\subsection{Structure}

The immediate format encodes a full 16-bit immediate value within the instruction.

The format has a single register field. For ALU operations that save a result it is both the first source operand and
the destination, so \texttt{ADDI r1, \#5} computes \reg{r1} = \reg{r1} + 5; there is no separate destination field.
Test operations read it as a source only, \mnemonic{MOVI} writes it without reading it, \mnemonic{LOAD} writes it, and
\mnemonic{STOR} reads it.

\textbf{Note:} Due to the large 16-bit operand, shift is not supported in this format. AF = 10 is still encodable; with no shift fields the shifter result is treated as zero flags, so it clears N, Z, C and V (Chapter~\ref{ch:alu-instructions}).

\begin{figure}[H]
	\centering
	\begin{bytefield}[bitwidth=1.3em]{32}
		\bitheader{0-31} \\
		\bitbox{6}{Opcode} &
		\bitbox{4}{Reg} &
		\bitbox{16}{Immediate Value} &
		\bitbox{2}{AF} &
		\bitbox{4}{Cond}
	\end{bytefield}
	\caption{Immediate Instruction Format}
	\label{fig:immediate-format}
\end{figure}

\begin{description}
	\item[Opcode (bits 31--26)] 6-bit operation code
	\item[Register (bits 25--22)] 4-bit register identifier; its role depends on the instruction (see above). For \mnemonic{LDEA}
	      and \mnemonic{LDEL} (opcodes \texttt{0x18}--\texttt{0x1F}), this field instead holds a 2-bit address-register-pair
	      index in bits 23--22 (00 = \reg{l}, 01 = \reg{a}, 10 = \reg{s}, 11 = \reg{p}) naming the destination pair, using the
	      same encoding as the AF field. Bits 25--24 are reserved for these opcodes and are written as zero; a non-zero value is
	      architecturally undefined.
	\item[Immediate (bits 21--6)] 16-bit signed or unsigned immediate value
	\item[Additional Flags (bits 5--4)] 2-bit address register pair selector (for memory operations)
	      \begin{itemize}
		      \item 00 = \reg{l} (link register)
		      \item 01 = \reg{a} (address register)
		      \item 10 = \reg{s} (stack pointer)
		      \item 11 = \reg{p} (program counter)
	      \end{itemize}
	      or 2-bit status register update source selector (for ALU operations)
	      \begin{itemize}
		      \item 00 = Status register not updated
		      \item 01 = Status register updated from ALU operation
		      \item 10 = Status register updated from shift operation
		      \item 11 = Reserved; N, Z, C and V are architecturally undefined (the operation still completes and writes its result), and assemblers never emit this encoding
	      \end{itemize}
	\item[Condition (bits 3--0)] 4-bit condition code
\end{description}

\subsection{Encoding Examples}

\begin{table}[H]
	\centering
	\small
	\begin{tabular}{lcccccc}
		\toprule
		\textbf{Instruction}          & \textbf{Opcode} & \textbf{Reg} & \textbf{Immediate}    & \textbf{AF} & \textbf{Cond} & \textbf{Hex}        \\
		\midrule
		\texttt{ADDI r1, \#100}       & 0x00            & 0x0          & 0x0064                & 01          & 0x0           & \texttt{0x00001910} \\
		                              & (000000)        & (0000)       & (0000 0000 0110 0100) & (01)        & (0000)        &                     \\
		\midrule
		\texttt{CMPI r2, \#0x8000}    & 0x0A            & 0x1          & 0x8000                & 01          & 0x0           & \texttt{0x28600010} \\
		                              & (001010)        & (0001)       & (1000 0000 0000 0000) & (01)        & (0000)        &                     \\
		\midrule
		\texttt{LOAD r3, (\#16, a)}   & 0x14            & 0x2          & 0x0010                & 01          & 0x0           & \texttt{0x50800410} \\
		                              & (010100)        & (0010)       & (0000 0000 0001 0000) & (01)        & (0000)        &                     \\
		\midrule
		\texttt{STOR (\#-2, s), r4}   & 0x10            & 0x3          & 0xFFFE                & 10          & 0x0           & \texttt{0x40FFFFA0} \\
		                              & (010000)        & (0011)       & (1111 1111 1111 1110) & (10)        & (0000)        &                     \\
		\midrule
		\texttt{LDEA a, (\#4, s)}     & 0x18            & 0x1          & 0x0004                & 10          & 0x0           & \texttt{0x60400120} \\
		                              & (011000)        & (0001)       & (0000 0000 0000 0100) & (10)        & (0000)        &                     \\
		\midrule
		\texttt{LDEA a, -(\#0, s)}    & 0x1A            & 0x1          & 0x0000                & 10          & 0x0           & \texttt{0x68400020} \\
		                              & (011010)        & (0001)       & (0000 0000 0000 0000) & (10)        & (0000)        &                     \\
		\midrule
		\texttt{LDEL p, (\#0, a)}     & 0x1C            & 0x3          & 0x0000                & 01          & 0x0           & \texttt{0x70C00010} \\
		                              & (011100)        & (0011)       & (0000 0000 0000 0000) & (01)        & (0000)        &                     \\
		\midrule
		\texttt{LDEL p, (\#0, a)+}    & 0x1E            & 0x3          & 0x0000                & 01          & 0x0           & \texttt{0x78C00010} \\
		                              & (011110)        & (0011)       & (0000 0000 0000 0000) & (01)        & (0000)        &                     \\
		\midrule
		\texttt{BRAN \#4}             & 0x18            & 0x3          & 0x0004                & 11          & 0x0           & \texttt{0x60C00130} \\
		                              & (011000)        & (0011)       & (0000 0000 0000 0100) & (11)        & (0000)        &                     \\
		\midrule
		\texttt{BRSR \#4}             & 0x1C            & 0x3          & 0x0004                & 11          & 0x0           & \texttt{0x70C00130} \\
		                              & (011100)        & (0011)       & (0000 0000 0000 0100) & (11)        & (0000)        &                     \\
		\midrule
		\texttt{ADDI|== r5, \#42}     & 0x00            & 0x4          & 0x002A                & 01          & 0x1           & \texttt{0x01000A91} \\
		                              & (000000)        & (0100)       & (0000 0000 0010 1010) & (01)        & (0001)        &                     \\
		\midrule
		\texttt{SUBI|>= r6, \#10}     & 0x02            & 0x5          & 0x000A                & 01          & 0xB           & \texttt{0x0940029B} \\
		                              & (000010)        & (0101)       & (0000 0000 0000 1010) & (01)        & (1011)        &                     \\
		\midrule
		\texttt{ADDI[N] r1, \#100}    & 0x00            & 0x0          & 0x0064                & 00          & 0x0           & \texttt{0x00001900} \\
		                              & (000000)        & (0000)       & (0000 0000 0110 0100) & (00)        & (0000)        &                     \\
		\midrule
		\texttt{CMPI[A] r2, \#0x8000} & 0x0A            & 0x1          & 0x8000                & 01          & 0x0           & \texttt{0x28600010} \\
		                              & (001010)        & (0001)       & (1000 0000 0000 0000) & (01)        & (0000)        &                     \\
		\bottomrule
	\end{tabular}
	\caption{Immediate Format Encoding Examples}
	\label{tab:immediate-format-examples}
\end{table}


\noindent\textbf{Notes:}
\begin{itemize}
	\item Condition code examples: \texttt{0x0} (0000) = Always, \texttt{0x1} (0001) = Equal (==), \texttt{0xB} (1011) = Signed
	      >=
	\item AF field for memory ops: 00=\reg{l}, 01=\reg{a}, 10=\reg{s}, 11=\reg{p}; for ALU ops: 00=no update, 01=update from ALU,
	      10=update from shift, 11=reserved (N, Z, C and V architecturally undefined; the result is still written)
	\item Negative immediates use two's complement (\imm{-2} = \texttt{0xFFFE})
\end{itemize}

\section{Short Immediate Format (with Shift)}

\subsection{Structure}

The short immediate format sacrifices immediate value width to include shift information, allowing for more powerful
single-instruction operations.

This format is supported only for ALU instructions; it cannot be used with branching or load/store instructions.

As in the immediate format, the single register field is both the first source operand and the destination for ALU
operations that save a result, so \texttt{ADDI r1, \#5, LSL \#2} computes \reg{r1} = (\reg{r1} shifted left by 2) + 5:
the shift applies to the register, not the immediate (Chapter~\ref{ch:shift-operations}). Test operations read it as a
source only.

\begin{figure}[H]
	\centering
	\begin{bytefield}[bitwidth=1.3em]{32}
		\bitheader{0-31} \\
		\bitbox{6}{Op} &
		\bitbox{4}{Reg} &
		\bitbox{8}{Imm} &
		\bitbox{1}{\scriptsize SO} &
		\bitbox{3}{ST} &
		\bitbox{4}{SA} &
		\bitbox{2}{AF} &
		\bitbox{4}{Cond}
	\end{bytefield}
	\caption{Short Immediate with Shift Format}
	\label{fig:short-imm-format}
\end{figure}

\begin{description}
	\item[Opcode (bits 31--26)] 6-bit operation code
	\item[Register (bits 25--22)] 4-bit register identifier; for ALU operations that save a result it is both the first source
	      and the destination (see above)
	\item[Immediate (bits 21--14)] 8-bit unsigned immediate value, zero-extended to 16 bits before use
	\item[Shift Operand (bit 13)] Determines how the shift amount field is interpreted
	      \begin{itemize}
		      \item 0 = Shift amount field contains the literal shift amount
		      \item 1 = Shift amount field refers to register that contains shift amount
	      \end{itemize}
	\item[Shift Type (bits 12--10)] Type of shift operation (see Table~\ref{tab:shift-types} for encoding details)
	\item[Shift Amount (bits 9--6)] Shift amount (or register containing shift amount)
	\item[Additional Flags (bits 5--4)] 2-bit status register update source selector
	      \begin{itemize}
		      \item 00 = Status register not updated
		      \item 01 = Status register updated from ALU operation
		      \item 10 = Status register updated from shift operation
		      \item 11 = Reserved; N, Z, C and V are architecturally undefined (the operation still completes and writes its result), and assemblers never emit this encoding
	      \end{itemize}
	\item[Condition (bits 3--0)] 4-bit condition code
\end{description}

\subsection{Encoding Examples}

\begin{table}[H]
	\centering
	\footnotesize
	\setlength{\tabcolsep}{2.5pt}
	\begin{tabular}{lccccccccc}
		\toprule
		\textbf{Instruction}               & \textbf{Op} & \textbf{Reg} & \textbf{Imm} & \textbf{SO} & \textbf{ST} & \textbf{SA} & \textbf{AF} & \textbf{Cond} & \textbf{Hex}        \\
		\midrule
		\texttt{ADDI r1, \#2, LSL \#3}     & 0x20        & 0x0          & 0x02         & 0           & 001         & 0x3         & 01          & 0x0           & \texttt{0x800084D0} \\
		                                   & (100000)    & (0000)       & (00000010)   & (0)         & (001)       & (0011)      & (01)        & (0000)        &                     \\
		\midrule
		\texttt{ORRI r2, \#1, LSL \#10}    & 0x25        & 0x1          & 0x01         & 0           & 001         & 0xA         & 01          & 0x0           & \texttt{0x94404690} \\
		                                   & (100101)    & (0001)       & (00000001)   & (0)         & (001)       & (1010)      & (01)        & (0000)        &                     \\
		\midrule
		\texttt{ANDI r3, \#7, ASR \#2}     & 0x24        & 0x2          & 0x07         & 0           & 100         & 0x2         & 01          & 0x0           & \texttt{0x9081D090} \\
		                                   & (100100)    & (0010)       & (00000111)   & (0)         & (100)       & (0010)      & (01)        & (0000)        &                     \\
		\midrule
		\texttt{SUBI r4, \#5, LSL r2}      & 0x22        & 0x3          & 0x05         & 1           & 001         & 0x1         & 01          & 0x0           & \texttt{0x88C16450} \\
		                                   & (100010)    & (0011)       & (00000101)   & (1)         & (001)       & (0001)      & (01)        & (0000)        &                     \\
		\midrule
		\texttt{XORI|!= r7, \#15, LSR \#4} & 0x26        & 0x6          & 0x0F         & 0           & 010         & 0x4         & 01          & 0x2           & \texttt{0x9983C912} \\
		                                   & (100110)    & (0110)       & (00001111)   & (0)         & (010)       & (0100)      & (01)        & (0010)        &                     \\
		\midrule
		\texttt{CMPI|NS r1, \#8, RTL r3}   & 0x2A        & 0x0          & 0x08         & 1           & 101         & 0x2         & 01          & 0x5           & \texttt{0xA8023495} \\
		                                   & (101010)    & (0000)       & (00001000)   & (1)         & (101)       & (0010)      & (01)        & (0101)        &                     \\
		\midrule
		\texttt{SHFT r3, RTR \#5}          & 0x25        & 0x2          & 0x00         & 0           & 110         & 0x5         & 10          & 0x0           & \texttt{0x94801960} \\
		                                   & (100101)    & (0010)       & (00000000)   & (0)         & (110)       & (0101)      & (10)        & (0000)        &                     \\
		\midrule
		\texttt{ADDI[N] r1, \#2, LSL \#3}  & 0x20        & 0x0          & 0x02         & 0           & 001         & 0x3         & 00          & 0x0           & \texttt{0x800084C0} \\
		                                   & (100000)    & (0000)       & (00000010)   & (0)         & (001)       & (0011)      & (00)        & (0000)        &                     \\
		\midrule
		\texttt{ORRI[S] r2, \#1, LSL \#10} & 0x25        & 0x1          & 0x01         & 0           & 001         & 0xA         & 10          & 0x0           & \texttt{0x944046A0} \\
		                                   & (100101)    & (0001)       & (00000001)   & (0)         & (001)       & (1010)      & (10)        & (0000)        &                     \\
		\bottomrule
	\end{tabular}
	\caption{Short Immediate Format Encoding Examples}
	\label{tab:short-imm-format-examples}
\end{table}


\noindent\textbf{Notes:}
\begin{itemize}
	\item SO (Shift Operand): 0 = literal shift amount, 1 = register contains shift amount (see rows 4 and 6)
	\item ST (Shift Type): see Table~\ref{tab:shift-encoding}
	\item AF field for ALU: 00=no update, 01=update from ALU, 10=update from shift (see SHFT row 7), 11=reserved (N, Z, C and V
	      architecturally undefined; the result is still written)
	\item Condition code examples: \texttt{0x0} = Always, \texttt{0x2} = Not Equal (!=), \texttt{0x5} = Negative Set (NS, N = 1)
	\item SHFT is a meta-instruction using ORRI opcode \opcode{25} with AF=10
\end{itemize}

\section{Register Format}

\subsection{Structure}

The register format allows operations on up to three registers with optional shift.

\begin{figure}[H]
	\centering
	\begin{bytefield}[bitwidth=1.3em]{32}
		\bitheader{0-31} \\
		\bitbox{6}{Op} &
		\bitbox{4}{R1} &
		\bitbox{4}{R2} &
		\bitbox{4}{R3} &
		\bitbox{1}{\scriptsize SO} &
		\bitbox{3}{ST} &
		\bitbox{4}{SA} &
		\bitbox{2}{AF} &
		\bitbox{4}{Cond}
	\end{bytefield}
	\caption{Register Instruction Format}
	\label{fig:register-format}
\end{figure}

\begin{description}
	\item[Opcode (bits 31--26)] 6-bit operation code
	\item[Register 1 (bits 25--22)] 4-bit destination register. For \mnemonic{LDEA} and \mnemonic{LDEL} (opcodes
	      \texttt{0x18}--\texttt{0x1F}), this field instead holds a 2-bit address-register-pair index in bits 23--22 (00 =
	      \reg{l}, 01 = \reg{a}, 10 = \reg{s}, 11 = \reg{p}) naming the destination pair, using the same encoding as the AF
	      field. Bits 25--24 are reserved for these opcodes and are written as zero; a non-zero value is architecturally
	      undefined.
	\item[Register 2 (bits 21--18)] 4-bit source A register
	\item[Register 3 (bits 17--14)] 4-bit source B register
	\item[Shift Operand (bit 13)] Determines how the shift amount field is interpreted
	      \begin{itemize}
		      \item 0 = Shift amount field contains the literal shift amount
		      \item 1 = Shift amount field refers to register that contains shift amount
	      \end{itemize}
	\item[Shift Type (bits 12--10)] Type of shift operation (see Table~\ref{tab:shift-types} for encoding details)
	\item[Shift Amount (bits 9--6)] Shift amount (or register containing shift amount)
	\item[Additional Flags (bits 5--4)] 2-bit address register pair selector (for memory operations)
	      \begin{itemize}
		      \item 00 = \reg{l} (link register)
		      \item 01 = \reg{a} (address register)
		      \item 10 = \reg{s} (stack pointer)
		      \item 11 = \reg{p} (program counter)
	      \end{itemize}
	      or 2-bit status register update source selector (for ALU operations)
	      \begin{itemize}
		      \item 00 = Status register not updated
		      \item 01 = Status register updated from ALU operation
		      \item 10 = Status register updated from shift operation
		      \item 11 = Reserved; N, Z, C and V are architecturally undefined (the operation still completes and writes its result), and assemblers never emit this encoding
	      \end{itemize}
	\item[Condition (bits 3--0)] 4-bit condition code
\end{description}

\subsection{Register Operand Semantics}

For most three-operand instructions:
\begin{itemize}
	\item R1 = Destination
	\item R2 = First source operand
	\item R3 = Second source operand (optional)
\end{itemize}

\textbf{Note:} If only two register operands are given, the assembler uses R1 as both destination and
first source operand when encoding the instruction:

\begin{lstlisting}
ADDR r1, r2, r3  ; r1 = r2 + r3
ADDR r1, r2      ; r1 = r1 + r2 (R2 implicitly = R1)
\end{lstlisting}

\subsection{Encoding Examples}

\begin{table}[H]
	\centering
	\footnotesize
	\setlength{\tabcolsep}{2.5pt}
	\begin{tabular}{lcccccccccc}
		\toprule
		\textbf{Instruction}                & \textbf{Op} & \textbf{R1} & \textbf{R2} & \textbf{R3} & \textbf{SO} & \textbf{ST} & \textbf{SA} & \textbf{AF} & \textbf{Cond} & \textbf{Hex}        \\
		\midrule
		\texttt{ADDR r1, r2, r3}            & 0x30        & 0x0         & 0x1         & 0x2         & 0           & 000         & 0x0         & 01          & 0x0           & \texttt{0xC0048010} \\
		                                    & (110000)    & (0000)      & (0001)      & (0010)      & (0)         & (000)       & (0000)      & (01)        & (0000)        &                     \\
		\midrule
		\texttt{ADDR r1, r2, r3, LSL \#2}   & 0x30        & 0x0         & 0x1         & 0x2         & 0           & 001         & 0x2         & 01          & 0x0           & \texttt{0xC0048490} \\
		                                    & (110000)    & (0000)      & (0001)      & (0010)      & (0)         & (001)       & (0010)      & (01)        & (0000)        &                     \\
		\midrule
		\texttt{LOAD r1, (r2, a)}           & 0x15        & 0x0         & 0x0         & 0x1         & 0           & 000         & 0x0         & 01          & 0x0           & \texttt{0x54004010} \\
		                                    & (010101)    & (0000)      & (0000)      & (0001)      & (0)         & (000)       & (0000)      & (01)        & (0000)        &                     \\
		\midrule
		\texttt{STOR -(r2, s), r1}          & 0x13        & 0x0         & 0x0         & 0x1         & 0           & 000         & 0x0         & 10          & 0x0           & \texttt{0x4C004020} \\
		                                    & (010011)    & (0000)      & (0000)      & (0001)      & (0)         & (000)       & (0000)      & (10)        & (0000)        &                     \\
		\midrule
		\texttt{LDEA a, (r2, s)}            & 0x19        & 0x1         & 0x0         & 0x1         & 0           & 000         & 0x0         & 10          & 0x0           & \texttt{0x64404020} \\
		                                    & (011001)    & (0001)      & (0000)      & (0001)      & (0)         & (000)       & (0000)      & (10)        & (0000)        &                     \\
		\midrule
		\texttt{LDEA a, -(r2, s)}           & 0x1B        & 0x1         & 0x0         & 0x1         & 0           & 000         & 0x0         & 10          & 0x0           & \texttt{0x6C404020} \\
		                                    & (011011)    & (0001)      & (0000)      & (0001)      & (0)         & (000)       & (0000)      & (10)        & (0000)        &                     \\
		\midrule
		\texttt{LDEL p, (r2, a)}            & 0x1D        & 0x3         & 0x0         & 0x1         & 0           & 000         & 0x0         & 01          & 0x0           & \texttt{0x74C04010} \\
		                                    & (011101)    & (0011)      & (0000)      & (0001)      & (0)         & (000)       & (0000)      & (01)        & (0000)        &                     \\
		\midrule
		\texttt{LDEL p, (r2, a)+}           & 0x1F        & 0x3         & 0x0         & 0x1         & 0           & 000         & 0x0         & 01          & 0x0           & \texttt{0x7CC04010} \\
		                                    & (011111)    & (0011)      & (0000)      & (0001)      & (0)         & (000)       & (0000)      & (01)        & (0000)        &                     \\
		\midrule
		\texttt{LJMP a, r2}                 & 0x19        & 0x3         & 0x0         & 0x1         & 0           & 000         & 0x0         & 01          & 0x0           & \texttt{0x64C04010} \\
		                                    & (011001)    & (0011)      & (0000)      & (0001)      & (0)         & (000)       & (0000)      & (01)        & (0000)        &                     \\
		\midrule
		\texttt{LJSR a, r2}                 & 0x1D        & 0x3         & 0x0         & 0x1         & 0           & 000         & 0x0         & 01          & 0x0           & \texttt{0x74C04010} \\
		                                    & (011101)    & (0011)      & (0000)      & (0001)      & (0)         & (000)       & (0000)      & (01)        & (0000)        &                     \\
		\midrule
		\texttt{CMPR r4, r5}                & 0x3A        & 0x3         & 0x3         & 0x4         & 0           & 000         & 0x0         & 01          & 0x0           & \texttt{0xE8CD0010} \\
		                                    & (111010)    & (0011)      & (0011)      & (0100)      & (0)         & (000)       & (0000)      & (01)        & (0000)        &                     \\
		\midrule
		\texttt{XORR r2, r3, r4, RTR \#1}   & 0x36        & 0x1         & 0x2         & 0x3         & 0           & 110         & 0x1         & 01          & 0x0           & \texttt{0xD848D850} \\
		                                    & (110110)    & (0001)      & (0010)      & (0011)      & (0)         & (110)       & (0001)      & (01)        & (0000)        &                     \\
		\midrule
		\texttt{SUBR|>> r1, r2, r3, ASR r4} & 0x32        & 0x0         & 0x1         & 0x2         & 1           & 100         & 0x3         & 01          & 0xD           & \texttt{0xC804B0DD} \\
		                                    & (110010)    & (0000)      & (0001)      & (0010)      & (1)         & (100)       & (0011)      & (01)        & (1101)        &                     \\
		\midrule
		\texttt{ANDR|CS r5, r6, r7, LSL r1} & 0x34        & 0x4         & 0x5         & 0x6         & 1           & 001         & 0x0         & 01          & 0x3           & \texttt{0xD115A413} \\
		                                    & (110100)    & (0100)      & (0101)      & (0110)      & (1)         & (001)       & (0000)      & (01)        & (0011)        &                     \\
		\midrule
		\texttt{ADDR[N] r1, r2, r3}         & 0x30        & 0x0         & 0x1         & 0x2         & 0           & 000         & 0x0         & 00          & 0x0           & \texttt{0xC0048000} \\
		                                    & (110000)    & (0000)      & (0001)      & (0010)      & (0)         & (000)       & (0000)      & (00)        & (0000)        &                     \\
		\midrule
		\texttt{SUBR[S] r1, r2, r3, ASR r4} & 0x32        & 0x0         & 0x1         & 0x2         & 1           & 100         & 0x3         & 10          & 0x0           & \texttt{0xC804B0E0} \\
		                                    & (110010)    & (0000)      & (0001)      & (0010)      & (1)         & (100)       & (0011)      & (10)        & (0000)        &                     \\
		\bottomrule
	\end{tabular}
	\caption{Register Format Encoding Examples}
	\label{tab:register-format-examples}
\end{table}


\noindent\textbf{Notes:}
\begin{itemize}
	\item When only two register operands are given, R2 is encoded as a copy of R1 and the second operand is encoded in R3, for
	      example \texttt{ADDR r4, r5} uses R1=\texttt{0x3}, R2=\texttt{0x3}, R3=\texttt{0x4}
	\item SO (Shift Operand): 0 = literal shift amount, 1 = register contains shift amount (see the \texttt{SUBR|>>} and
	      \texttt{ANDR|CS} rows)
	\item ST (Shift Type): see Table~\ref{tab:shift-encoding}
	\item AF for memory ops: 00=\reg{l}, 01=\reg{a}, 10=\reg{s}, 11=\reg{p}; for ALU ops: 00=no update, 01=update from ALU,
	      10=update from shift, 11=reserved (N, Z, C and V architecturally undefined; the result is still written)
	\item Condition code examples: \texttt{0x0} = Always, \texttt{0xD} = Signed > (>>), \texttt{0x3} = Carry Set (CS)
\end{itemize}

\section{Operand Encoding Details}

\subsection{Condition Codes}

The condition field occupies bits 3--0 of every instruction. The sixteen condition-code encodings and their meanings
are tabulated once, in Table~\ref{tab:condition-codes} of Chapter~\ref{ch:condition-codes}, which is the authoritative
copy.

\subsection{Shift Types}

The shift-type field occupies bits 12--10 of the formats that carry one (see the field descriptions above). The
shift-type encodings are tabulated once, in Table~\ref{tab:shift-encoding} of Chapter~\ref{ch:shift-operations}, which
is the authoritative copy.

\section{Opcode Encoding Patterns}

The SIRCIS instruction set uses systematic opcode patterns to simplify decoding:

\subsection{Opcode Groups}

\begin{itemize}
	\item \textbf{\texttt{0x0\_}}: Immediate operand
	\item \textbf{\texttt{0x1\_}}: Memory operations
	\item \textbf{\texttt{0x2\_}}: Short immediate with shift
	\item \textbf{\texttt{0x3\_}}: Register operand
\end{itemize}

\subsection{Test vs. Save}

This pattern applies to the ALU groups \texttt{0x0\_}, \texttt{0x2\_}, and \texttt{0x3\_} only. Each ALU operation has
both a "save result" and a "test only" variant:

\begin{itemize}
	\item \textbf{\texttt{0x\_0}--\texttt{0x\_6}}: Save result to destination
	\item \textbf{\texttt{0x\_8}--\texttt{0x\_E}}: Test only (update flags, do not save the result)
	\item \textbf{\texttt{0x\_F}}: Coprocessor call (\mnemonic{COPI}/\mnemonic{COPR}), not a test-only
	      form
\end{itemize}

Adding \texttt{0x8} to a "save result" opcode in the range \texttt{0x\_0}--\texttt{0x\_6} gives its "test only" form.
Opcode \texttt{0x\_7} (\mnemonic{MOVI}/\mnemonic{MOVR}) has no test-only form. Opcodes \texttt{0x27} and \texttt{0x2F}
are exceptions to this pattern: they are undocumented opcodes (see Appendix~\ref{appendix:undocumented}), not a "save
result" operation or a coprocessor call.

\begin{table}[H]
	\centering
	\begin{tabular}{clll}
		\toprule
		\textbf{Operation}  & \textbf{Save result}              & \textbf{Test only}                & \textbf{Test-only opcodes (Imm / Short / Reg)} \\
		\midrule
		Add                 & \mnemonic{ADDI} / \mnemonic{ADDR} & \mnemonic{CMNI} / \mnemonic{CMNR} & \opcode{08} / \opcode{28} / \opcode{38}        \\
		Add with carry      & \mnemonic{ADCI} / \mnemonic{ADCR} & \mnemonic{CNCI} / \mnemonic{CNCR} & \opcode{09} / \opcode{29} / \opcode{39}        \\
		Subtract            & \mnemonic{SUBI} / \mnemonic{SUBR} & \mnemonic{CMPI} / \mnemonic{CMPR} & \opcode{0A} / \opcode{2A} / \opcode{3A}        \\
		Subtract with carry & \mnemonic{SBCI} / \mnemonic{SBCR} & \mnemonic{CPCI} / \mnemonic{CPCR} & \opcode{0B} / \opcode{2B} / \opcode{3B}        \\
		AND                 & \mnemonic{ANDI} / \mnemonic{ANDR} & \mnemonic{TSAI} / \mnemonic{TSAR} & \opcode{0C} / \opcode{2C} / \opcode{3C}        \\
		OR                  & \mnemonic{ORRI} / \mnemonic{ORRR} & \mnemonic{TSOI} / \mnemonic{TSOR} & \opcode{0D} / \opcode{2D} / \opcode{3D}        \\
		XOR                 & \mnemonic{XORI} / \mnemonic{XORR} & \mnemonic{TSXI} / \mnemonic{TSXR} & \opcode{0E} / \opcode{2E} / \opcode{3E}        \\
		\bottomrule
	\end{tabular}
	\caption{Save and Test-Only ALU Instructions}
	\label{tab:save-and-test-instructions}
\end{table}

For example:

\begin{itemize}
	\item \opcode{04}: \mnemonic{ANDI} -- AND immediate, save result
	\item \opcode{0C}: \mnemonic{TSAI} -- Test AND immediate (does not save the result)
\end{itemize}

\section{Decode Trees}
\label{sec:decode-trees}

The opcode alone determines how the rest of a 32-bit instruction is decoded. The figures in this section show that
process as a decision tree for each opcode group. In every figure the top strip is the instruction word, drawn to
scale. The shaded boxes are the opcode bits consulted at each level of the tree, with the same shading on the matching
tree node, and the unshaded boxes are operand fields, whose interpretation is given in the box below the tree.
\texttt{op[n]} means opcode bit n, which is instruction bit n + 26. Appendix~\ref{appendix:opcode-map} lists every
opcode in a single table.

% ---- Decode tree helpers ------------------------------------------------------------------------
\newcommand{\dtu}{0.46}% width of one instruction bit in the word strip, in cm

\tikzset{
	dt node/.style={draw, rounded corners=1.5pt, align=center, font=\footnotesize, inner sep=3pt},
	dt edge/.style={-{Latex[length=1.8mm]}, thick},
	dt label/.style={font=\scriptsize\ttfamily, inner sep=1.5pt}
}

% \dtbox{first bit column}{width in bits}{fill}{label}{instruction bit range}
\newcommand{\dtbox}[5]{%
	\draw[fill=#3] ({\dtu*(#1)},0) rectangle ({\dtu*((#1)+(#2))},-0.55);
	\node[font=\footnotesize\ttfamily] at ({\dtu*((#1)+(#2)/2)},-0.275) {#4};
	\node[font=\tiny, text=sircgray] at ({\dtu*((#1)+(#2)/2)},-0.78) {#5};
}

% \dtrow{op[2:0]}{opcode}{mnemonic}{description}
\newcommand{\dtrow}[4]{\texttt{#1} & \opcode{#2} & \mnemonic{#3} & #4\\}
\newcommand{\dtrowx}[4]{\textcolor{sircgray}{\texttt{#1}} & \textcolor{sircgray}{\texttt{0x#2}} & \textcolor{sircgray}{#3} & \textcolor{sircgray}{#4}\\}

% Save rows 0x_0--0x_6 and the move at 0x_7; #1 = opcode high digit, #2 = mnemonic suffix
\newcommand{\dtsaverows}[2]{%
	\dtrow{000}{#10}{ADD#2}{Add}
	\dtrow{001}{#11}{ADC#2}{Add with carry}
	\dtrow{010}{#12}{SUB#2}{Subtract}
	\dtrow{011}{#13}{SBC#2}{Subtract with carry}
	\dtrow{100}{#14}{AND#2}{AND}
	\dtrow{101}{#15}{ORR#2}{OR}
	\dtrow{110}{#16}{XOR#2}{XOR}%
}
% Test rows 0x_8--0x_E
\newcommand{\dttestrows}[2]{%
	\dtrow{000}{#18}{CMN#2}{Compare negative}
	\dtrow{001}{#19}{CNC#2}{Compare negative with carry}
	\dtrow{010}{#1A}{CMP#2}{Compare}
	\dtrow{011}{#1B}{CPC#2}{Compare with carry}
	\dtrow{100}{#1C}{TSA#2}{Test AND}
	\dtrow{101}{#1D}{TSO#2}{Test OR}
	\dtrow{110}{#1E}{TSX#2}{Test XOR}%
}

% \dtalu{format node text}{strip boxes after the opcode}{op[3]=0 op[2:0]=111 row}{op[3]=1 op[2:0]=111 row}{ALU suffix}{opcode high digit}{legend rows}
\newcommand{\dtalu}[7]{%
	\begin{tikzpicture}
		% Instruction word strip: the opcode is split into the three parts that drive the tree
		\dtbox{0}{2}{black!4}{5:4}{31--30}
		\dtbox{2}{1}{black!10}{3}{29}
		\dtbox{3}{3}{black!16}{2:0}{28--26}
		#2
		% Tree
		\node[dt node, fill=black!4, text width=2.5cm] (fmt) at (1.45,-5.4) {#1};
		\node[dt node, fill=black!10] (b3) at (5.0,-5.4) {\texttt{op[3]}};
		\node[dt node, fill=black!16, anchor=north west, inner sep=4pt] (save) at (7.6,-1.7)
		{\begin{tabular}{@{}c@{\hspace{6pt}}c@{\hspace{6pt}}l@{\hspace{6pt}}l@{}}
				\multicolumn{4}{@{}l@{}}{\textbf{Save result:} \texttt{op[2:0]} selects the operation} \\[2pt]
				\dtsaverows{#6}{#5}
				#3
			\end{tabular}};
		\node[dt node, fill=black!16, anchor=north west, inner sep=4pt] (test) at (7.6,-6.1)
		{\begin{tabular}{@{}c@{\hspace{6pt}}c@{\hspace{6pt}}l@{\hspace{6pt}}l@{}}
				\multicolumn{4}{@{}l@{}}{\textbf{Test only:} \texttt{op[2:0]} selects the operation} \\[2pt]
				\dttestrows{#6}{#5}
				#4
			\end{tabular}};
		\draw[dt edge] (fmt) -- (b3);
		\draw[dt edge] (b3.east) -- ++(0.5,0) |- (save.west) node[dt label, pos=0.82, above] {0};
		\draw[dt edge] (b3.east) -- ++(0.5,0) |- (test.west) node[dt label, pos=0.82, above] {1};
		% Operand fields
		\node[draw, rounded corners=1.5pt, anchor=north west, inner sep=4pt, font=\footnotesize] at (0,-10.5)
		{\begin{tabular}{@{}p{2.9cm}p{11.4cm}@{}}
				\multicolumn{2}{@{}l@{}}{\textbf{Operand fields} (read the same way for every opcode above, unless noted)} \\[2pt]
				#7
			\end{tabular}};
	\end{tikzpicture}%
}

% ---- Figures --------------------------------------------------------------------------------------

\begin{figure}[H]
	\centering
	\renewcommand{\arraystretch}{1}
	\dtalu{\texttt{op[5:4]\,=\,00}\\Immediate format}
	{%
		\dtbox{6}{4}{white}{Reg}{25--22}
		\dtbox{10}{16}{white}{Immediate}{21--6}
		\dtbox{26}{2}{white}{AF}{5--4}
		\dtbox{28}{4}{white}{Cond}{3--0}
	}
	{\dtrow{111}{07}{MOVI}{Move}}
	{\dtrow{111}{0F}{COPI$^{*}$}{Coprocessor call}}
	{I}{0}
	{%
		Reg (25--22) & \texttt{rD}, the destination, for save forms; \texttt{rS}, the register tested, for test-only forms\\
		Immediate (21--6) & 16-bit value. This format has no shift fields\\
		AF (5--4) & Status update source: \texttt{00} none, \texttt{01} ALU result, \texttt{10} shifter result, \texttt{11} reserved (N, Z, C and V architecturally undefined)\\
		Cond (3--0) & One of 16 condition codes (Table~\ref{tab:condition-code-encoding})\\[2pt]
		\multicolumn{2}{@{}l@{}}{\scriptsize $^{*}$\mnemonic{COPI} is a coprocessor call, not a test form; its fields are described in Chapter~\ref{ch:coprocessor}.}\\
	}
	\caption{Decode tree for the immediate ALU group (\texttt{0x0\_})}
	\label{fig:decode-tree-immediate}
\end{figure}

\begin{figure}[H]
	\centering
	\renewcommand{\arraystretch}{1}
	\dtalu{\texttt{op[5:4]\,=\,10}\\Short immediate\\format}
	{%
		\dtbox{6}{4}{white}{Reg}{25--22}
		\dtbox{10}{8}{white}{Imm}{21--14}
		\dtbox{18}{1}{white}{SO}{13}
		\dtbox{19}{3}{white}{ST}{12--10}
		\dtbox{22}{4}{white}{SA}{9--6}
		\dtbox{26}{2}{white}{AF}{5--4}
		\dtbox{28}{4}{white}{Cond}{3--0}
	}
	{\dtrowx{111}{27}{--}{Undocumented, Appendix~\ref{appendix:undocumented}}}
	{\dtrowx{111}{2F}{--}{Undocumented, Appendix~\ref{appendix:undocumented}}}
	{I}{2}
	{%
		Reg (25--22) & \texttt{rD}, the destination, for save forms; \texttt{rS}, the register tested, for test-only forms\\
		Imm (21--14) & 8-bit unsigned value, zero-extended to 16 bits\\
		SO, ST, SA (13--6) & \texttt{ST} selects the shift type (Table~\ref{tab:shift-types}). If \texttt{SO} is 0, \texttt{SA} is the literal shift count; if 1, \texttt{SA} names the register holding the count. The shift applies to the register operand, never the immediate\\
		AF (5--4) & Status update source: \texttt{00} none, \texttt{01} ALU result, \texttt{10} shifter result, \texttt{11} reserved (N, Z, C and V architecturally undefined)\\
		Cond (3--0) & One of 16 condition codes (Table~\ref{tab:condition-code-encoding})\\
	}
	\caption{Decode tree for the short immediate ALU group (\texttt{0x2\_})}
	\label{fig:decode-tree-short-immediate}
\end{figure}

\begin{figure}[H]
	\centering
	\renewcommand{\arraystretch}{1}
	\dtalu{\texttt{op[5:4]\,=\,11}\\Register format}
	{%
		\dtbox{6}{4}{white}{R1}{25--22}
		\dtbox{10}{4}{white}{R2}{21--18}
		\dtbox{14}{4}{white}{R3}{17--14}
		\dtbox{18}{1}{white}{SO}{13}
		\dtbox{19}{3}{white}{ST}{12--10}
		\dtbox{22}{4}{white}{SA}{9--6}
		\dtbox{26}{2}{white}{AF}{5--4}
		\dtbox{28}{4}{white}{Cond}{3--0}
	}
	{\dtrow{111}{37}{MOVR}{Move}}
	{\dtrow{111}{3F}{COPR$^{*}$}{Coprocessor call}}
	{R}{3}
	{%
		R1 (25--22) & \texttt{rD}, the destination, for save forms; ignored by test-only forms\\
		R2 (21--18) & \texttt{rS1}, the left operand. \mnemonic{MOVR} does not use it\\
		R3 (17--14) & \texttt{rS2}, the right operand (for \mnemonic{MOVR}, the source)\\
		SO, ST, SA (13--6) & \texttt{ST} selects the shift type (Table~\ref{tab:shift-types}). If \texttt{SO} is 0, \texttt{SA} is the literal shift count; if 1, \texttt{SA} names the register holding the count. The shift applies to \texttt{rS1}\\
		AF (5--4) & Status update source: \texttt{00} none, \texttt{01} ALU result, \texttt{10} shifter result, \texttt{11} reserved (N, Z, C and V architecturally undefined)\\
		Cond (3--0) & One of 16 condition codes (Table~\ref{tab:condition-code-encoding})\\[2pt]
		\multicolumn{2}{@{}l@{}}{\scriptsize $^{*}$\mnemonic{COPR} is a coprocessor call, not a test form; its fields are described in Chapter~\ref{ch:coprocessor}.}\\
	}
	\caption{Decode tree for the register ALU group (\texttt{0x3\_})}
	\label{fig:decode-tree-register}
\end{figure}

% Memory and control group (0x1_)
% \dtmem{y}{op1 node text}{op[1]=0 leaf}{op[1]=1 leaf}{label for op[3:2]}
\newcommand{\dtmemfamily}[6]{%
	% #1 index (0--3), #2 op[3:2] bits, #3 family mnemonic, #4 leaf 0 text, #5 leaf 1 text, #6 opcode high/low base digit
	\pgfmathsetmacro{\ya}{-2.45-(2*#1)*1.05}
	\pgfmathsetmacro{\yb}{-2.45-(2*#1+1)*1.05}
	\pgfmathsetmacro{\ym}{(\ya+\yb)/2}
	\node[dt node, fill=black!16] (fam#1) at (7.9,\ym) {\mnemonic{#3}\\\texttt{op[1]}};
	\node[dt node, fill=black!24, anchor=west, text width=4.7cm] (leafa#1) at (9.5,\ya) {#4};
	\node[dt node, fill=black!24, anchor=west, text width=4.7cm] (leafb#1) at (9.5,\yb) {#5};
	\draw[dt edge] (op32.east) -- ++(0.45,0) |- (fam#1.west) node[dt label, pos=0.8, above] {#2};
	\draw[dt edge] (fam#1.east) -- ++(0.3,0) |- (leafa#1.west) node[dt label, pos=0.8, above] {0};
	\draw[dt edge] (fam#1.east) -- ++(0.3,0) |- (leafb#1.west) node[dt label, pos=0.8, above] {1};
}

\begin{figure}[H]
	\centering
	\renewcommand{\arraystretch}{1}
	\begin{tikzpicture}
		\dtbox{0}{2}{black!4}{5:4}{31--30}
		\dtbox{2}{2}{black!10}{3:2}{29--28}
		\dtbox{4}{1}{black!16}{1}{27}
		\dtbox{5}{1}{black!24}{0}{26}
		\dtbox{6}{26}{white}{Operand fields (layout chosen by \texttt{op[0]})}{25--0}
		\node[dt node, fill=black!4, text width=2.5cm] (fmt) at (1.45,-5.7) {\texttt{op[5:4]\,=\,01}\\Memory and control};
		\node[dt node, fill=black!10] (op32) at (4.9,-5.7) {\texttt{op[3:2]}};
		\draw[dt edge] (fmt) -- (op32);
		\node[font=\scriptsize, anchor=south west, align=left, text width=5.1cm, inner sep=1pt] at (9.6,-1.8) {\texttt{op[0]} then picks the format: 0 = \mbox{Immediate}, 1 = \mbox{Register}};
		\dtmemfamily{0}{00}{STOR}{\textbf{Indirect}\\\opcode{10} imm \,/\, \opcode{11} reg}{\textbf{Pre-decrement}\\\opcode{12} imm \,/\, \opcode{13} reg}{}
		\dtmemfamily{1}{01}{LOAD}{\textbf{Indirect}\\\opcode{14} imm \,/\, \opcode{15} reg}{\textbf{Post-increment}\\\opcode{16} imm \,/\, \opcode{17} reg}{}
		\dtmemfamily{2}{10}{LDEA}{\textbf{Indirect}\\\opcode{18} imm \,/\, \opcode{19} reg}{\textbf{Pre-decrement}\\\opcode{1A} imm \,/\, \opcode{1B} reg}{}
		\dtmemfamily{3}{11}{LDEL}{\textbf{Indirect}\\\opcode{1C} imm \,/\, \opcode{1D} reg}{\textbf{Post-increment}\\\opcode{1E} imm \,/\, \opcode{1F} reg}{}
		\node[draw, rounded corners=1.5pt, anchor=north west, inner sep=4pt, font=\footnotesize] at (0,-11.4)
		{\begin{tabular}{@{}p{1.9cm}p{6.1cm}p{6.3cm}@{}}
				\multicolumn{3}{@{}l@{}}{\textbf{Operand fields} (chosen by \texttt{op[0]})}                                                                                                                                                                                                                                                                                                                     \\[2pt]
				          & \textbf{Immediate} (\texttt{op[0]} = 0)                                                                                                                                    & \textbf{Register} (\texttt{op[0]} = 1)                                                                                                                                                                  \\[2pt]
				25--22    & \texttt{rD} for \mnemonic{LOAD}, \texttt{rS} for \mnemonic{STOR}. \mnemonic{LDEA} and \mnemonic{LDEL} put the destination pair in bits 23--22 (25--24 are zero)            & As for the immediate format, except that \mnemonic{STOR} leaves them unused                                                                                                                             \\
				21--6     & 16-bit displacement added to the address register pair                                                                                                                     & \texttt{rO}, the offset register (17--14). \mnemonic{STOR} also uses \texttt{rS} (21--18) and a shift of \texttt{rS} (\texttt{SO}, \texttt{ST}, \texttt{SA}); those fields are zero for the other three \\
				5--4 (AF) & \multicolumn{2}{@{}p{12.4cm}@{}}{Address register pair used for the effective address: \texttt{00} \reg{l}, \texttt{01} \reg{a}, \texttt{10} \reg{s}, \texttt{11} \reg{p}}                                                                                                                                                                                                           \\
				3--0      & \multicolumn{2}{@{}p{12.4cm}@{}}{Condition code}                                                                                                                                                                                                                                                                                                                                     \\
			\end{tabular}};
	\end{tikzpicture}
	\caption{Decode tree for the memory and control group (\texttt{0x1\_})}
	\label{fig:decode-tree-memory}
\end{figure}

\section{Instruction Fetch and Alignment}

\subsection{Fetch Process}

Since instructions are 32 bits and the data bus is 16 bits wide, each instruction requires two fetch cycles:

\begin{enumerate}
	\item Fetch the high word from the address in \reg{pl}.
	\item Fetch the low word from \reg{pl} + 1.
	\item After decode, \reg{pl} advances by 2.
\end{enumerate}

The address of the instruction is latched before \reg{pl} advances. Every read of \reg{p} and every address calculation
that selects \reg{p} uses that latched address, so they see the instruction's own address even though \reg{pl} has
already advanced by the time the address is calculated in phase 3. A link value is the latched address plus 2, the
address of the next instruction.

\subsection{Alignment Requirements}

Instructions span two consecutive words and must begin at an even word address. The program counter (\reg{pl}) must
always be even.

Fetching an instruction from an odd word address raises an alignment fault.

\section{Format Selection Guidelines}

When writing assembly code, the assembler automatically selects the appropriate format based on operands:

\begin{table}[H]
	\centering
	\small
	\begin{tabular}{lll}
		\toprule
		\textbf{Assembly}               & \textbf{Format Used} & \textbf{Notes}          \\
		\midrule
		\texttt{ADDI r1, \#1000}        & Immediate            & Full 16-bit immediate   \\
		\texttt{ADDI r1, \#10, LSL \#2} & Short Immediate      & 8-bit immediate + shift \\
		\texttt{ADDR r1, r2, r3}        & Register             & Three register operands \\
		\bottomrule
	\end{tabular}
	\caption{Format Selection Examples}
	\label{tab:format-selection}
\end{table}

The choice of format affects instruction encoding but is generally transparent to the programmer, as the assembler
handles the details.
